% Anti-windup on one step, and the rate mismatch drawn to scale.
\begin{tikzpicture}[font=\sffamily\scriptsize, x=1mm, y=1mm]

% ================================================================ the step
\node[chlabel, anchor=west] at (0,46) {one setpoint step into a saturating output, three times};
\draw[taxis] (0,10) -- (124,10);
\draw[taxis] (0,10) -- (0,42);
\draw[deadline] (0,32) -- (124,32);
\node[tlabel, anchor=west] at (2,34) {the setpoint};

% no anti-windup: large overshoot, slow return
\draw[signal, draw=red!70!black, line width=0.9pt]
  (4,10) .. controls (18,10) and (26,42) .. (40,42)
  -- (62,42) .. controls (78,42) and (84,32) .. (104,32);
\node[tlabel, anchor=west, text=red!60!black] at (64,44) {none: it winds up, and pays for it};

% clamping: modest overshoot
\draw[signal, draw=black, line width=0.9pt]
  (4,10) .. controls (18,10) and (24,36) .. (36,36)
  .. controls (46,36) and (50,32) .. (66,32);
\node[tlabel, anchor=west] at (38,37.5) {clamping};

% back-calculation: smooth
\draw[signal, draw=blue!55!black, line width=0.9pt]
  (4,10) .. controls (18,10) and (24,32.5) .. (34,32.2) -- (66,32);
\node[tlabel, anchor=west] at (36,29) {back-calculation, written for this chapter};

\draw[tspan] (4,6) -- node[tlabel]{saturated for this long} (26,6);

% ================================================================ the rate mismatch
\node[chlabel, anchor=west] at (0,-2) {and the rate mismatch, drawn to scale};
\node[signame] at (-2,-10) {the host};
\draw[taxis] (0,-13) -- (124,-13);
\foreach \x in {0,40,80,120}{
  \draw[release] (\x,-17) -- (\x,-8);
  \node[fieldbox, fill=extcol, minimum width=12mm, text width=10mm, minimum height=4.5mm,
        font=\sffamily\tiny] at ({\x+1},-12.5) {setpoint};
}
\node[tlabel, anchor=west] at (0,-19.5) {100 Hz, the framework's own default, on a kernel with no real-time guarantee};

\node[signame] at (-2,-27) {the node};
\draw[taxis] (0,-30) -- (124,-30);
\foreach \x in {0,4,...,120}{ \draw[release] (\x,-33) -- (\x,-27); }
\node[tlabel, anchor=west] at (2,-24.5) {1 kHz: ten corrections between one host update and the next};

% ================================================================ following error
\node[chlabel, anchor=west] at (0,-38) {and the quantity the joint is judged on};
\node[regcell, fill=usrcol, minimum width=124mm, minimum height=7mm, anchor=south west,
      inner sep=0pt] at (0,-53.5) {};
\draw[taxis] (0,-50) -- (124,-50);
\node[tlabel, anchor=west] at (2,-46) {the tolerance band};
\draw[signal, line width=0.9pt]
  (0,-50) .. controls (10,-44) and (18,-56) .. (30,-50)
  .. controls (42,-45) and (52,-54) .. (66,-50)
  .. controls (78,-38) and (86,-40) .. (96,-50) -- (124,-50);
\draw[faultedge, -] (76,-36) -- (92,-36);
\node[tlabel, anchor=west, text=red!60!black] at (94,-36) {outside: chapter 15's code -4};

\node[umlnote, text width=60mm, anchor=north west] at (0,-60)
  {The three traces above are the chapter's own measurement, run on the same
   saturating plant with the same gains and the same step. The comparison does
   not exist in the open literature for a small embedded controller, which is
   why producing it is worth more than another tuning guide.};
\node[umlnote, text width=60mm, anchor=north west] at (64,-60)
  {The band is not decoration. When the error leaves it the node reports the
   path tolerance code chapter 15 borrowed, which a controller on the host
   already knows how to act on, instead of a private number that would have to
   be explained to every consumer.};
\end{tikzpicture}
